\section{Mô hình miền nghiệp vụ}

Thiết kế chi tiết tập trung vào các đối tượng nghiệp vụ có ảnh hưởng trực tiếp đến tính đúng
đắn của hệ thống: sản phẩm, SKU, số dư tồn, sổ cái tồn kho, đơn hàng, dòng đơn hàng, giữ hàng
và ảnh chụp giá vốn.

Các dịch vụ ứng dụng điều phối ca sử dụng nhưng không chứa những quy tắc dữ liệu cốt lõi thuộc
miền nghiệp vụ. Điều này giúp quy tắc giữ hàng và tính giá vốn có thể được kiểm thử độc lập.

\section{Thiết kế sổ cái tồn kho}

Khi phát sinh biến động tồn kho, quy trình tổng quát là:

\begin{enumerate}
    \item bắt đầu giao dịch cơ sở dữ liệu;
    \item xác định SKU và chi nhánh;
    \item khóa bản ghi số dư cần thay đổi;
    \item kiểm tra điều kiện nghiệp vụ;
    \item cập nhật số dư;
    \item ghi một bản ghi mới vào sổ cái;
    \item ghi tham chiếu đến nghiệp vụ nguồn;
    \item xác nhận giao dịch.
\end{enumerate}

Không xóa hoặc sửa các bản ghi lịch sử của sổ cái trong quy trình thông thường.

\section{Máy trạng thái đơn hàng}

Trạng thái trung tâm của đơn hàng là:

\[
Draft \rightarrow Reserved \rightarrow Confirmed \rightarrow Completed
\]

Trạng thái \textit{Cancelled} có thể được sử dụng từ các trạng thái trung gian hợp lệ. Việc
chuyển trạng thái phải được kiểm tra tại lớp nghiệp vụ để tránh chuyển trực tiếp giữa các trạng
thái không được phép.

\section{Tính giá vốn bình quân gia quyền}

Sau khi nhập thêm hàng, giá vốn bình quân được cập nhật theo công thức:

\[
C_{WAC} =
\frac{Q_{old}C_{old} + Q_{in}C_{in}}
     {Q_{old}+Q_{in}}
\]

Khi bán $Q_{sold}$ đơn vị, giá vốn của giao dịch là:

\[
COGS = Q_{sold} \times C_{WAC}
\]

Giá vốn được chụp lại tại thời điểm hoàn tất đơn hàng để báo cáo lịch sử không thay đổi khi
giá vốn bình quân của SKU thay đổi trong tương lai.

\section{Thuật toán giữ hàng}

Quy trình giữ hàng gồm:

\begin{enumerate}
    \item bắt đầu giao dịch;
    \item tải bản ghi tồn của SKU và chi nhánh;
    \item khóa bản ghi số dư;
    \item tính $Available = OnHand - Reserved$;
    \item từ chối nếu số lượng yêu cầu lớn hơn $Available$;
    \item nếu đủ, tăng $Reserved$;
    \item tạo bản ghi giữ hàng;
    \item xác nhận giao dịch.
\end{enumerate}

Khóa được giữ trong khoảng thời gian cần thiết để ngăn các giao dịch cạnh tranh cùng sử dụng
một số lượng tồn.

\section{Giải phóng và trừ tồn}

Khi đơn hàng bị hủy sau khi đã giữ hàng, số lượng giữ phải được giảm và trạng thái giữ được
đóng. Khi đơn hàng hoàn tất, số lượng thực tế được trừ khỏi tồn và lượng giữ tương ứng được
loại khỏi số lượng đang giữ. Các thay đổi đều phải có bản ghi lịch sử phù hợp.

\section{Thanh toán nhanh tại POS}

Quy trình thanh toán POS gồm tìm SKU, tạo giỏ, xác nhận số lượng, giữ hàng, xác nhận đơn,
trừ tồn, lưu COGS và trả dữ liệu biên nhận. Chuỗi Giữ $\rightarrow$ Xác nhận $\rightarrow$
Trừ tồn phải được thiết kế nguyên tử đối với phần dữ liệu có ảnh hưởng đến tồn kho.

\section{Xử lý lỗi và giao dịch bù trừ}

Nếu một bước trong giao dịch thất bại, giao dịch phải được hủy để tránh trạng thái dở dang.
Đối với các tác động bên ngoài không thể hoàn tác bằng một giao dịch cơ sở dữ liệu duy nhất,
hệ thống sử dụng giao dịch bù trừ thay vì xóa lịch sử.

\section{Xử lý đồng thời}

Với các SKU có tần suất bán cao, có thể ưu tiên khóa bi quan trên bản ghi số dư. Các thao tác
không xung đột có thể sử dụng cơ chế kiểm tra phiên bản phù hợp. Cách triển khai cuối cùng phải
được đánh giá trong Chương~\ref{c:evaluation}.
